(a) Every homogenization of an element of $I(Y)$ vanishes on $Y$, hence on $\overline Y$. Conversely, if homogeneous $F$ vanishes on $\overline Y$, its dehomogenization $F(1,x_1,\ldots,x_n)$ lies in $I(Y)$, and $F$ is a power of $x_0$ times its homogenization. Thus $I(\overline Y)$ is generated by $\beta(I(Y))$.

(b) In affine coordinates, $I(Y)=(y-x^2,z-x^3)$ by \exref{I.1.2}. In projective coordinates $(u,x,y,z)$,
\[I(\overline Y)=(uy-x^2,uz-xy,xz-y^2).\]
These relations reduce every monomial to one of $u^az^d,u^axz^d,u^ayz^d$. For fixed total degree, their images under $(u,x,y,z)\mapsto(s^3,s^2t,st^2,t^3)$ are distinct monomials, since their $t$-exponents are $3d,3d+1,3d+2$, respectively. Thus the relations generate the kernel, which is $I(\overline Y)$. The homogenizations $uy-x^2,u^2z-x^3$ do not generate this ideal: their degree-two part is spanned by $uy-x^2$, and does not contain $uz-xy$.
